\documentclass[12pt]{article}

\usepackage[margin=1in]{geometry}
\usepackage{amsmath,amssymb}
\usepackage{booktabs}
\usepackage{listings}
\usepackage{xcolor}
\usepackage{hyperref}
\usepackage{setspace}
\usepackage{longtable}

\onehalfspacing

\lstdefinestyle{sql}{
  language=SQL,
  basicstyle=\small\ttfamily,
  keywordstyle=\bfseries\color{blue!70!black},
  commentstyle=\color{gray},
  stringstyle=\color{red!60!black},
  showstringspaces=false,
  breaklines=true,
  frame=single,
  backgroundcolor=\color{gray!5},
  morekeywords={ENUM,FLOAT,TEXT,BOOLEAN,REFERENCES,CHECK}
}

\title{Supplementary Information:\\Schema for Fuzzy-Joined Data Products}
\author{Gaurav Sood}
\date{\today}

\begin{document}
\maketitle

\section*{Overview}

This supplement proposes a concrete relational schema for shipping data products built on fuzzy joins. The schema covers the two-table case: a source table $S$ (e.g., firms) linked to a target table $T$ (e.g., villages) via string similarity on noisy identifiers. Five tables are defined. The first two are inputs; the remaining three are artifacts the data producer ships alongside the analysis-ready panel.

All schemas use SQL-style definitions for precision. Primary keys are marked \texttt{PK}; foreign keys reference their parent table. Implementers can adapt these to flat files (CSV/Parquet) by treating the key constraints as documentation rather than enforcement.

\section{Input Tables}

\subsection{Source Table}

The source table contains one row per record to be linked. In the running example, these are firms from an Economic Census. The table includes the raw reported name used for matching, outcome variables, and any covariates measured at the source level.

\begin{lstlisting}[style=sql]
CREATE TABLE source (
    source_id       INTEGER PRIMARY KEY,
    reported_name   TEXT NOT NULL,
    -- outcome and source-level covariates
    revenue         FLOAT,
    n_employees     INTEGER,
    year            INTEGER
);
\end{lstlisting}

\subsection{Target Table}

The target table contains one row per entity in the reference frame. In the running example, these are census villages. The table includes the canonical name, geographic identifiers, and target-level covariates used in downstream regressions.

\begin{lstlisting}[style=sql]
CREATE TABLE target (
    target_id       INTEGER PRIMARY KEY,
    target_name     TEXT NOT NULL,
    -- geographic hierarchy
    district_id     INTEGER,
    state_id        INTEGER,
    -- target-level covariates
    literacy        FLOAT,
    population      INTEGER,
    area_sq_km      FLOAT
);
\end{lstlisting}

\section{Shipped Artifacts}

\subsection{Canonical Linkage Table}

One row per source record. This is the default analysis dataset: the deterministic resolution of the candidate-match graph under a stated policy $(\tau, \delta)$. Every source record appears exactly once, whether matched or not.

\begin{lstlisting}[style=sql]
CREATE TABLE canonical_linkage (
    source_id       INTEGER PRIMARY KEY
                    REFERENCES source(source_id),
    target_id       INTEGER
                    REFERENCES target(target_id),
                    -- NULL if unmatched
    match_score     FLOAT,
                    -- score of the assigned pair; NULL if unmatched
    match_status    ENUM('exact',
                         'unambiguous_fuzzy',
                         'ambiguous_dropped',
                         'below_threshold',
                         'no_candidates')
                    NOT NULL,
    match_method    TEXT,
                    -- e.g., 'levenshtein', 'masala_merge_v2'
    score_gap       FLOAT,
                    -- s_{i(1)} - s_{i(2)}; NULL if <= 1 candidate
    n_candidates    INTEGER NOT NULL DEFAULT 0,
    component_id    INTEGER
                    -- connected component in the candidate graph
                    -- (see Section 2.4)
);
\end{lstlisting}

\paragraph{Design notes.}
\begin{itemize}
    \item \texttt{match\_status} distinguishes the reason a record is or is not linked. \texttt{exact} means the match was deterministic (e.g., a shared unique identifier). \texttt{unambiguous\_fuzzy} means the top candidate passed both $\tau$ and $\delta$. The remaining three categories are different failure modes: ambiguous (multiple close candidates), below threshold (best score too low), and no candidates (nothing passed screening).
    \item \texttt{target\_id} is \texttt{NULL} for unmatched records. This makes the canonical linkage a left join from source to target, not an inner join.  Every source record is present, so the analyst can see what was dropped and why.
    \item \texttt{component\_id} identifies the connected component in the candidate graph (Section~2.4). Records with a unique candidate have a singleton component. Records in confusable clusters share a component. This lets the analyst quickly identify which records are affected by many-to-many ambiguity.
    \item \texttt{n\_candidates} and \texttt{score\_gap} are summary statistics from the candidate set. They allow the analyst to assess match quality at the row level without loading the full candidate table.
\end{itemize}

\paragraph{Example rows.}

\begin{center}
\small
\begin{tabular}{rrrccrr}
\toprule
source\_id & target\_id & score & status & gap & n\_cand & comp \\
\midrule
0 & 0 & 1.00 & unambig\_fuzzy & 0.14 & 3 & 1 \\
3 & \texttt{NULL} & \texttt{NULL} & ambig\_dropped & 0.00 & 2 & 1 \\
6 & 1 & 1.00 & unambig\_fuzzy & 0.14 & 3 & 1 \\
11 & \texttt{NULL} & \texttt{NULL} & ambig\_dropped & 0.00 & 2 & 1 \\
30 & 5 & 1.00 & exact & 1.00 & 1 & 12 \\
\bottomrule
\end{tabular}
\end{center}

Firm~0 (rampur, reported ``rampur'') is matched to target~0 (rampur) with score 1.00 and gap 0.14 over the second candidate. Firm~3 (rampur, reported ``rampure'') is dropped because the gap is zero. Firm~6 (rampur, reported ``rampura'') is a false positive matched to target~1 (rampura). All four are in component~1 (the ram-cluster). Firm~30 (dharamkot) is an exact match with a single candidate, in its own component.

\subsection{Candidate-Match Table}

One row per (source, candidate target) pair above the screening threshold $\varepsilon$. This is the bipartite candidate-match graph $\mathcal{G}$. Every resolution policy discussed in the paper can be reconstructed from this table.

\begin{lstlisting}[style=sql]
CREATE TABLE candidate_matches (
    source_id       INTEGER NOT NULL
                    REFERENCES source(source_id),
    target_id       INTEGER NOT NULL
                    REFERENCES target(target_id),
    match_score     FLOAT NOT NULL
                    CHECK (match_score >= 0 AND match_score <= 1),
    match_rank      INTEGER NOT NULL,
                    -- rank within this source_id's candidate set
                    -- (1 = best score)
    is_default      BOOLEAN NOT NULL,
                    -- TRUE if this pair is the canonical assignment
    component_id    INTEGER,
                    -- same component_id as in canonical_linkage
    PRIMARY KEY (source_id, target_id)
);
\end{lstlisting}

\paragraph{Design notes.}
\begin{itemize}
    \item The primary key is (source\_id, target\_id): each pair appears at most once.
    \item \texttt{match\_rank} is the within-source ranking by score (1 = highest). Ties are broken arbitrarily but consistently.
    \item \texttt{is\_default} flags which pair, if any, is the canonical assignment. For unmatched source records, no row has \texttt{is\_default = TRUE}. For matched records, exactly one row does. This makes it trivial to reconstruct the canonical linkage: \texttt{SELECT * FROM candidate\_matches WHERE is\_default = TRUE}.
    \item Source records with zero candidates (status \texttt{no\_candidates} in the canonical linkage) have no rows in this table. To recover the full source population, left-join \texttt{canonical\_linkage} to \texttt{candidate\_matches}.
    \item The screening threshold $\varepsilon$ should be documented in metadata. It is typically looser than $\tau$ (the assignment threshold), so the candidate table includes pairs that the canonical policy would not assign.
\end{itemize}

\paragraph{Example rows.}

\begin{center}
\small
\begin{tabular}{rrlcrr}
\toprule
source\_id & target\_id & target\_name & score & rank & is\_default \\
\midrule
0 & 0 & rampur & 1.00 & 1 & TRUE \\
0 & 1 & rampura & 0.86 & 2 & FALSE \\
0 & 2 & ramnagar & 0.50 & 3 & FALSE \\
3 & 0 & rampur & 0.86 & 1 & FALSE \\
3 & 1 & rampura & 0.86 & 2 & FALSE \\
6 & 1 & rampura & 1.00 & 1 & TRUE \\
6 & 0 & rampur & 0.86 & 2 & FALSE \\
30 & 5 & dharamkot & 1.00 & 1 & TRUE \\
\bottomrule
\end{tabular}
\end{center}

Firm~3 has two candidate rows, neither flagged as default (it was dropped by the canonical policy). Firm~0 and Firm~6 each have a default row. Firm~30 has a single candidate.

\subsection{Validation Sample Table}

One row per source record in the verified golden set. The golden set is a random sample from the \emph{full source population}, including records that the canonical policy dropped. For each record, the true target is verified by a human or by a unique identifier available for a subset.

\begin{lstlisting}[style=sql]
CREATE TABLE validation_sample (
    source_id       INTEGER PRIMARY KEY
                    REFERENCES source(source_id),
    true_target_id  INTEGER NOT NULL
                    REFERENCES target(target_id),
    verification    ENUM('unique_id', 'clerical_review',
                         'field_verification')
                    NOT NULL,
    -- the canonical assignment for comparison
    canonical_target_id  INTEGER
                    REFERENCES target(target_id),
                    -- NULL if this record was dropped
    match_status    ENUM('exact',
                         'unambiguous_fuzzy',
                         'ambiguous_dropped',
                         'below_threshold',
                         'no_candidates')
                    NOT NULL
);
\end{lstlisting}

\paragraph{Design notes.}
\begin{itemize}
    \item The critical design requirement, discussed in the main text, is that this table must sample from the full source population, not just from matched records. Without dropped records, selection bias from false negatives cannot be assessed.
    \item \texttt{canonical\_target\_id} is the assignment from the canonical linkage (NULL if dropped). Including it alongside \texttt{true\_target\_id} lets the analyst compute the confusion matrix and calibration slope directly: regress $X_{\texttt{true\_target\_id}}$ on $X_{\texttt{canonical\_target\_id}}$ among matched records.
    \item \texttt{verification} records how the true link was established. This matters for assessing the reliability of the golden set itself.
    \item The golden set should be stratified to oversample from ambiguous connected components (high \texttt{n\_candidates}, low \texttt{score\_gap} in the canonical linkage). If stratified, sampling weights should be documented so the analyst can reweight to population representativeness.
\end{itemize}

\paragraph{Example rows.}

\begin{center}
\small
\begin{tabular}{rrrll}
\toprule
source\_id & true\_target & canonical\_target & status & verification \\
\midrule
0 & 0 & 0 & unambig\_fuzzy & clerical \\
3 & 0 & \texttt{NULL} & ambig\_dropped & clerical \\
6 & 0 & 1 & unambig\_fuzzy & clerical \\
30 & 5 & 5 & exact & unique\_id \\
\bottomrule
\end{tabular}
\end{center}

Firm~0: correctly matched (true = canonical = rampur). Firm~3: dropped by canonical policy, but verified as rampur. Firm~6: false positive (true = rampur, canonical = rampura). Firm~30: exact match confirmed by unique identifier.

\subsection{Connected Component Table}

One row per connected component in the candidate-match graph. A connected component is a maximal set of source and target records linked by candidate edges. Components with a single source and single target are singletons (no ambiguity). Dense components with multiple sources and targets are where matching decisions are interdependent.

\begin{lstlisting}[style=sql]
CREATE TABLE components (
    component_id    INTEGER PRIMARY KEY,
    n_sources       INTEGER NOT NULL,
    n_targets       INTEGER NOT NULL,
    n_edges         INTEGER NOT NULL,
    -- summary diagnostics
    density         FLOAT,
                    -- n_edges / (n_sources * n_targets)
    mean_score      FLOAT,
    min_gap         FLOAT,
                    -- minimum score gap among source records
    n_ambiguous     INTEGER,
                    -- sources with status 'ambiguous_dropped'
    x_range         FLOAT
                    -- max(X) - min(X) across targets in component
                    -- measures how much MI draws can vary
);
\end{lstlisting}

\paragraph{Design notes.}
\begin{itemize}
    \item This table is a diagnostic summary. The analyst can quickly identify components where matching uncertainty concentrates: large \texttt{n\_ambiguous}, high \texttt{x\_range}, high \texttt{density}.
    \item \texttt{x\_range} is computed for a specific covariate $X$ from the target table. If multiple covariates are of interest, this column should be replicated or the table extended.
    \item The independent-draw MI procedure in the main text treats each source record independently. For components with high density and large \texttt{x\_range}, the analyst may want to inspect draws at the component level to check whether independent draws produce implausible village-level firm counts.
\end{itemize}

\paragraph{Example rows.}

\begin{center}
\small
\begin{tabular}{rrrrcccc}
\toprule
comp\_id & sources & targets & edges & density & mean\_$s$ & n\_ambig & $X$ range \\
\midrule
1 & 21 & 3 & 60 & 0.95 & 0.72 & 4 & 0.35 \\
2 & 9 & 2 & 15 & 0.83 & 0.78 & 1 & 0.20 \\
3 & 1 & 1 & 1 & 1.00 & 1.00 & 0 & 0.00 \\
\bottomrule
\end{tabular}
\end{center}

Component~1 is the ram-cluster: 21 firms, 3 villages, nearly fully connected, 4 ambiguous firms, covariate range of 0.35 (literacy from 0.50 to 0.85). Component~2 is the sultan-cluster. Component~3 is a singleton (distinctive village).

\section{Reconstruction Queries}

The following queries reconstruct the resolution policies discussed in the main text from the shipped tables.

\paragraph{Canonical linkage (analysis-ready panel).}

\begin{lstlisting}[style=sql]
SELECT s.source_id, s.revenue, t.literacy
FROM source s
JOIN canonical_linkage cl ON s.source_id = cl.source_id
JOIN target t ON cl.target_id = t.target_id
WHERE cl.match_status IN ('exact', 'unambiguous_fuzzy');
\end{lstlisting}

\paragraph{Canonical linkage under alternative $(\tau', \delta')$.}

\begin{lstlisting}[style=sql]
-- Reconstruct canonical linkage with stricter thresholds
WITH ranked AS (
    SELECT source_id, target_id, match_score, match_rank,
           match_score - LEAD(match_score) OVER (
               PARTITION BY source_id ORDER BY match_rank
           ) AS gap
    FROM candidate_matches
)
SELECT source_id, target_id, match_score
FROM ranked
WHERE match_rank = 1
  AND match_score >= 0.70   -- tau'
  AND (gap >= 0.15 OR gap IS NULL);  -- delta'
\end{lstlisting}

\paragraph{Expanded full join.}

\begin{lstlisting}[style=sql]
-- One row per candidate pair, score-proportional weights
WITH totals AS (
    SELECT source_id, SUM(match_score) AS total_score
    FROM candidate_matches
    WHERE match_score >= 0.50  -- tau
    GROUP BY source_id
)
SELECT cm.source_id, cm.target_id, s.revenue, t.literacy,
       cm.match_score / tot.total_score AS weight
FROM candidate_matches cm
JOIN source s ON cm.source_id = s.source_id
JOIN target t ON cm.target_id = t.target_id
JOIN totals tot ON cm.source_id = tot.source_id
WHERE cm.match_score >= 0.50;
\end{lstlisting}

\paragraph{Collapsed (score-weighted average $X$).}

\begin{lstlisting}[style=sql]
-- One row per source, X = score-weighted average literacy
WITH totals AS (
    SELECT source_id, SUM(match_score) AS total_score
    FROM candidate_matches
    WHERE match_score >= 0.50
    GROUP BY source_id
)
SELECT cm.source_id, s.revenue,
       SUM(cm.match_score / tot.total_score * t.literacy)
           AS literacy_avg
FROM candidate_matches cm
JOIN source s ON cm.source_id = s.source_id
JOIN target t ON cm.target_id = t.target_id
JOIN totals tot ON cm.source_id = tot.source_id
WHERE cm.match_score >= 0.50
GROUP BY cm.source_id, s.revenue;
\end{lstlisting}

\paragraph{MI draw (single imputation).}
MI draws require randomness and are better implemented in a scripting language (Python, R, Stata). The pseudocode for one draw:

\begin{lstlisting}[style=sql]
-- Pseudocode: one MI draw
FOR EACH source_id IN canonical_linkage:
  IF match_status IN ('exact', 'unambiguous_fuzzy'):
    drawn_target = target_id  -- deterministic
  ELSE IF n_candidates > 0:
    candidates = SELECT target_id, match_score
                 FROM candidate_matches
                 WHERE source_id = :source_id
                   AND match_score >= :tau
    -- draw one target proportional to match_score
    drawn_target = WEIGHTED_SAMPLE(candidates)
  ELSE:
    SKIP  -- no candidates; exclude from this draw
  INSERT INTO mi_draw_m (source_id, drawn_target)
\end{lstlisting}

\paragraph{Regression calibration from the golden set.}

\begin{lstlisting}[style=sql]
-- Calibration slope: regress true X on assigned X
-- (matched records in golden set)
SELECT v.source_id,
       t_true.literacy AS x_true,
       t_canon.literacy AS x_assigned
FROM validation_sample v
JOIN target t_true ON v.true_target_id = t_true.target_id
JOIN target t_canon ON v.canonical_target_id = t_canon.target_id
WHERE v.canonical_target_id IS NOT NULL;
-- Feed into: x_true ~ x_assigned to get gamma_hat

-- Selection bias diagnostic (requires full-population golden set)
-- Oracle on full golden set:
SELECT v.source_id, s.revenue, t_true.literacy AS x_true
FROM validation_sample v
JOIN source s ON v.source_id = s.source_id
JOIN target t_true ON v.true_target_id = t_true.target_id;
-- Feed into: revenue ~ x_true (all records)
-- versus:    revenue ~ x_true (WHERE canonical_target IS NOT NULL)
-- Gap = selection bias from dropping ambiguous records
\end{lstlisting}


\section{Metadata}

The data product should ship a metadata file (JSON or plain text) documenting:

\begin{itemize}
    \item \textbf{Matching algorithm:} name, version, parameters (e.g., ``masala-merge v2.1, fuzziness = 0.8'').
    \item \textbf{Screening threshold} $\varepsilon$: the minimum score for a pair to appear in \texttt{candidate\_matches}.
    \item \textbf{Assignment thresholds} $(\tau, \delta)$: the score and gap thresholds used to produce the canonical linkage.
    \item \textbf{Score interpretation:} what the match score represents (e.g., ``1 minus normalized Levenshtein distance'', ``model-predicted match probability''), and whether raw scores should be treated as calibrated probabilities.
    \item \textbf{Validation sample design:} sampling frame (full population or matched only), stratification variables, sampling weights if applicable.
    \item \textbf{Coverage summary:} number of source records by match status, overall match rate, number of connected components, and the share of source records in non-singleton components.
\end{itemize}


\end{document}
